\chapter{Environment Setup and Sample Data}
\index{setup}\index{Fabric trial}\index{tools!installation}\index{sample data}\index{workspace conventions}%
This appendix is the standing setup reference for every lab in the book. Complete it once
before Chapter~0 and return whenever a lab assumes a tool, a connection, or a dataset that
your machine or tenant does not yet have. Everything here targets Windows with PowerShell,
matching the environment used throughout the program.

\section{Accounts and Tenant}
\index{Fabric trial!setup}%
\begin{itemize}
    \item \textbf{Fabric trial capacity.} Create a free 60-day trial capacity from the Fabric portal
    and enable Fabric in your tenant \cite{microsoftFabricTrial}. The trial covers every lab in the
    book; it expires, so it is unsuitable for anything with an SLA.
    \item \textbf{Tenant settings.} Where you administer the tenant, review the Fabric admin portal's
    tenant switches (Git integration, export, sharing) early---several Part~VI labs depend on them
    \cite{microsoftFabricTenantSettings}.
    \item \textbf{Azure access.} Some labs (gateways, Key Vault, Azure SQL source, Dataverse) need an
    Azure subscription or instructor-provided equivalents. A free Azure account suffices for the
    gateway and SQL labs.
\end{itemize}

\section{Local Tooling}
\index{VS Code}\index{Azure CLI}\index{SSMS}\index{Git!installation}%
Install once, in this order:

\begin{enumerate}
    \item \textbf{Python 3.x} --- used by the Week~13 producer, the Week~16 simulator, and synthetic
    data generation. Verify with \texttt{python --version}.
    \item \textbf{Git} --- required from Chapter~0 onward; Week~23 depends on it. Verify with
    \texttt{git --version}.
    \item \textbf{VS Code} --- with the Python extension; the editor for producer scripts and
    repository work.
    \item \textbf{Azure CLI (\texttt{az})} --- sign in with \texttt{az login}; used for Key Vault,
    storage, and Appendix~B's reference commands.
    \item \textbf{SQL Server Management Studio (SSMS)} --- the T-SQL client for the SQL analytics
    endpoint and warehouse labs (Chapters~5--6, 21).
    \item \textbf{Docker Desktop} (optional) --- the cleanest way to run the mock on-premises
    SQL Server for the Chapter~11 gateway lab.
    \item \textbf{Python packages} --- \texttt{azure-eventhub} (Week~13), \texttt{faker} (Weeks~16,
    24), \texttt{mlflow} and \texttt{scikit-learn} (Weeks~18--20).
\end{enumerate}

\begin{lstlisting}[language=bash,caption={One-time tool verification (PowerShell).},label={lst:appendixa-verify}]
python --version
git --version
az --version
az login
pip install azure-eventhub faker mlflow scikit-learn
\end{lstlisting}

\section{Repository and Credential Conventions}
\index{credentials!conventions}%
Clone or create the course repository with the standard layout---\texttt{data/}, \texttt{src/},
\texttt{tests/}, \texttt{infra/}, \texttt{docs/}---and observe three credential rules from day one:

\begin{itemize}
    \item Secrets live in environment variables (for example \texttt{AZURE\_CLIENT\_ID},
    \texttt{EVENTSTREAM\_CONN\_STR}) or Azure Key Vault---never in notebooks, pipeline JSON, or Git.
    \item Enable secret scanning on the course repository before Week~23's Git integration makes every
    workspace notebook a committed file.
    \item Commit generator scripts, never generated data; the datasets below are reproducible by design.
\end{itemize}

\section{Workspace Naming Conventions}
\index{workspace!naming}%
Labs reference workspaces by predictable names: \texttt{fabric-de-weekNN} for weekly labs,
\texttt{fabric-de-milestoneN} for capstones, and \texttt{fabric-de-w23-dev|test|prod} for the CI/CD
estate. Lakehouses are \texttt{lh\_<purpose>}, warehouses \texttt{dw\_<domain>}, eventstreams
\texttt{es\_<stream>}, KQL databases \texttt{kqldb\_<domain>}, pipelines \texttt{PL\_<action>}.
Consistency matters because runbooks, alerts, and the run ledger all reference these names.

\section{Sample Datasets}
\index{sample data!NYC Taxi}\index{sample data!Olist}\index{faker}%
The book's labs and capstones standardize on three data sources:

\begin{table}[H]
\centering
\small
\begin{tabular}{@{}p{3.2cm}p{4.6cm}p{5.8cm}@{}}
\toprule
\textbf{Dataset} & \textbf{Used in} & \textbf{Notes} \\
\midrule
NYC Taxi (TLC trip records) & Parts I--II (Chapters 1--8) & Public Parquet, mirrored on Azure Open Datasets; 1--5M rows suffice for capstones \\
Olist Brazilian e-commerce & Parts III and V (Chapters 9--12, 17--20) & $\sim$100k orders across 9 related CSVs: orders, items, customers, products, payments, reviews \\
Synthetic streams (\texttt{faker}) & Parts IV and VI (Chapters 13--16, 21--24) & Credit-card swipes, IoT telemetry, PII-style CRM rows; seed generators for reproducibility \\
\bottomrule
\end{tabular}
\caption{Standard datasets and where they are used.}
\label{tab:appendixa-datasets}
\end{table}

Volume guidance: Part~I medallion work at 1--10M CSV rows; the Part~II warehouse star schema at
$\sim$5M fact rows; Part~III re-ingests the nine Olist tables daily; Part~IV streams at
5--20k events per minute during demos; Part~V scores $\sim$100k rows per batch.

\begin{lstlisting}[language=Python,caption={Seeded synthetic PII generator (commit the script, not the data).},label={lst:appendixa-faker}]
from faker import Faker
import csv

fake = Faker()
Faker.seed(42)  # reproducible

with open("data/synthetic_customers.csv", "w", newline="") as f:
    w = csv.writer(f)
    w.writerow(["customer_id", "name", "email", "phone", "region"])
    for i in range(100_000):
        w.writerow([f"C{i:06d}", fake.name(), fake.email(),
                    fake.phone_number(), fake.random_element(["US", "EU", "APAC"])])
\end{lstlisting}

\section{Troubleshooting Checklist}
\index{troubleshooting!setup}%
\begin{itemize}
    \item \textbf{Notebook cannot see the lakehouse}: attach a default lakehouse to the notebook; verify
    the workspace is on a Fabric capacity, not a shared-only assignment.
    \item \textbf{\texttt{COPY INTO} path errors}: re-copy the OneLake DFS path from the lakehouse
    explorer; a wrong workspace or item segment is the usual cause (Chapter~5).
    \item \textbf{SSMS cannot connect}: use the SQL analytics endpoint connection string from lakehouse
    settings, Microsoft Entra MFA authentication, and confirm the endpoint reflects the table
    (files under \texttt{Files/} are invisible until registered as Delta tables).
    \item \textbf{Gateway offline}: check the gateway service on the host, then the cluster view in the
    portal; the Chapter~11 lab's failure drill covers the recovery pattern.
    \item \textbf{Throttling during labs}: the trial capacity is shared; pause unused capacities and
    check the Capacity Metrics app before assuming a lab is broken \cite{microsoftFabricMetricsApp}.
\end{itemize}

\begin{tipbox}
When a lab fails, reproduce the failure at the smallest possible scope---one row, one file, one
minute of telemetry---before investigating the estate. Every chapter's validation checklist is
written to support exactly that reduction.
\end{tipbox}
